\scanpage{50}
\mspage{II}

Soit $Y_1$ le normalisé de $Y$ en $y$, \inserted{\unsure{pt dés. point mt|pt dés. point} \unsure{chlt|clôt.}} \deleted{\unsure{Alors}}, en \unsure{le pt|tt} $y\in Y$. \unsure{Quitte} à se \deleted{localiser} \unsure{en} un pt maximal de la fibre \deleted{\unsure{spéciale}} \inserted{de $X_y$}, on \deleted{\unsure{prouve}} peut supposer $X_y$ \deleted{\unsure{fibre spéciale}} \unsure{irréduct.} Donc il en est \deleted{$X$ étant lisse sur $Y$, $X_1=X\times_Y Y_1$ l'est sur $Y_1$, $X_1$ est \unsure{normal}, donc \unsure{aussi} et} \inserted{de même de $X_1=X\times_Y Y_1$ \unsure{en $Y_1$|sur $Y_1$}, d'ailleurs le corps \unsure{résid.} \unsure{demeure}} \inserted{\unsure{celui} de $X_1$, et} le même. Il \unsure{coïncide} sur la fibre
\diagram{\begin{tikzcd}[ampersand replacement=\&,row sep=small,column sep=small]
X\arrow[dd] \& {}\\
{} \& X_1\arrow[ul]\arrow[dd]\\
Y \& {}\\
{} \& Y_1\arrow[ul]
\end{tikzcd}}
générique avec $Z_y$, il s'ensuit qu'on a un $X$-morphisme $X_1\mathrel{\deleted{\unsure{\to|\leftarrow}}}\leftarrow Z$ (qui est un isom sur les \unsure{points} de $(X_1)_y$). \deleted{\unsure{Comme} $Z$ \unsure{obtenu}, $X_1$ est \unsure{fidèle}, \unsure{puisque} sur $X$, $Y_1$ est \unsure{étale} sur $Y$.} \inserted{Par descente fid. plate, il \deleted{\unsure{faut}} s'ensuit que \deleted{\unsure{$X$}}$X_1$ est étale sur $Y$ en les pts de $(X_1)_y$, donc partout.} Donc on a
$Z\simeq X\times_Y Y_1$, $Y_1$ \unsure{revêtmt|revêtemt} fini étale de $Y$.

Par suite (revenant au cas général), \deleted{\unsure{$X_1$} est} $X_1$ est étale sur $X$ partout, et \deleted{\unsure{la comparaison}} \inserted{des considérations} de profondeur nous donnent un $X$-isom $X_1\simeq Z$, ce qui achève la démonstration.

\uline{NB} \unsure{On} a prouvé plus qu'on n'avait annoncé, savoir ceci : \unsure{On} a une suite exacte \deleted{\unsure{de}}
\margin{\unsure{TBE|TB} $\left.\rule{0pt}{3em}\right]$}
\diagram{\begin{tikzcd}[ampersand replacement=\&,column sep=small]
\pi_1(X_{\bar y},\unsure{\xi|\bar\xi})\arrow[r] \& \pi_1(X,\unsure{\xi|\zeta})\arrow[r] \& \pi_1(Y,\unsure{\zeta|\xi})\arrow[r] \& 0
\end{tikzcd}}

\unsure{Cela|Ceci} généralise (\unsure{modulo|en dehors de} l'hyp. noethérienne) SGA $\overline{\mathrm{IX}}$ 6.1., 6.2. \ldots. Ne pas confondre avec les résultats de SGA $\overline{\mathrm{X}}$, où on prend une fibre \uline{pas \inserted{\uline{néc}} générique} (en revanche, \unsure{3|?} il y faut supposer que $f$ est propre \ldots.).
% archive sheet number: 49
